\section{Introduction} \label{sec:introduction}

Modern Humanoids are able to achieve highly dynamic motions that their statically stable predecessors could only dream of \cite{tong_2024}.
Unitrees H1 and G2 execute impressive tasks, dancing \cite{sun_2026} and catching footballs like goalkeepers \cite{ren_2026}.
Having such powerful platforms at hand enables research towards extended evasion maneuvers.
These motions were previously not feasible due to the lack of power, stability, and joint range.
Yang et al. 2026 \cite{yang_2026} display such evasion by implementing a full dodgeball controller onto the G2.
\\
When it comes to handling collisions with approaching objects, previous robots, like the HRP-2 \cite{kaneko_2004} were mostly limited to impact resistance or slow evasion maneuvers
The field standing of push recovery is traditionally either dominated by "ankle" or "hip" strategies. \cite{hofmann_2006}
\todo{necessary?}
Hofmann \cite{hofmann_2006} mentions Allum and Honegger \cite{allum_1992}, as well as Nashner and McCollum \cite{nashner_1985}, as the foundational works in human push recovery studies.
Those strategies, however, are only feasible for small impact forces.
Higher forces, which exceed a certain threshold require capture steps \cite{pratt_2006}.
\\
In scenarios with low impact force compliance control can be used.
Compliance control describes methods to control a robot, which has to smooth out motions under the influence of external disturbances.
Such disturbances can be produced by impacts or error tolerances.
Per definition push resistance and locomotion tasks are different in there models, as the former requires an additional contact.
Hung et al. model those as a spring-damper system. \cite{huang_2022}
\\
However the superior solution would be to evade the approaching object and return to the initial pose without any contact.
This way the humanoid can calculate beforehand the resulting motion as well as counteract the possible stability difficulties.
Full collision avoidance also removes the risk of damage to both the humanoid and the object to evade.
As a result of there human-like statue, humanoids are for example set for applications in industry, health care or social service \cite{tong_2024,mccartney_2020}.
Such settings require the humanoid to be aware of other actors in its environment, like robotic arms or human coworkers.
Especially when working around humans it is essential to keep a reasonable safety distance.
This topic leads us to the well known research area of collision avoidance.

\paragraph{Collision Avoidance and Path Planning}
To solve the NP-hard problem of avoiding collisions for robotics \cite{canny_1988} different approaches have been developed.
A main distinction is the time frame in which the evasion motion has to be computed.
Works like \cite{yang_2026,kohlbrenner_2026} focus on real-time planning, with a closed loop and black-boxed system, spanning object detection, motion synthesis and control.
Recent systems of that kind are mostly driven by reinforcement or imitation learning providing little to no inside into there internals. \todo{bold statement?}
\\
The most important element remains generating the evasion motion.
By usage of a unified system for detection, path planning and motion execution, well understood components are reimplemented.
For the use case of evading approaching objects with humanoid robots already existing methods can be utilized.
An example is \cite{casas_2025}, which estimates object trajectories by fusing multiple sensors to a clean forecast.
Gab et al. 2022 \cite{gad_2022} also shows a variety of different \acrfull{mto} approaches.
\\
To finally execute the planned evasion motion a \acrlong{wbc} is required.
It stabilizes the robot and not only stabilizes the robot, but also counteracts environmental uncertainties and planning inaccuracies introduced by simplified models \cite{gu_2026}.
We showed that both the object tracking and motion execution fields offer a large variety of different methods.
Moving forward, we therefore focus on the generation of feasible evasion trajectories.
\\ % end of why we need only path planning
While path planning is a huge field in robotics only a small portion of the works are focusing on humanoids, which are evading approaching objects.
A lot of work is run on quadrupeds or an \acrfull{uav}. \cite{zhao_2026,zhang_2026,jacquet_2025}
In general, these pose a lower \acrfull{dof} problem and provide better stability.
However, as described before, some situations require robots to work in spaces specifically designed for humans.
Some even on robotic arms, which eliminates all stability concerns and further reduces the problem size \cite{mercy_2016,haviland_2021}.
These low \acrshort{dof} robots are utilize mostly model-based controllers, as there computation is fast enough for \acrfull{mpc}.
\\
When it comes to humanoids, recent work in the field of approaching object avoidance, is solved with learning-based methods \cite{yang_2026,kohlbrenner_2026}.
Solutions like Yang et al. \cite{yang_2026} display impressive results on open spaces, taking steps and moving out of the way when necessary.
These capture step resembling motions are less taxing on the dynamic stability limits of the humanoid.
Environments in industrial or health care settings can require tightly packed spaces, in which evasive maneuvers should be concealed to anchored balancing motions.
Rapid and large stepping motions in contrast, possibly threaten a collision with the robots surroundings, or in worst case other actors in the scene.
Therefore stepping is not always an option humanoids, when recovering from pushes or avoiding those impacts directly \cite{huang_2022}.
Most of the current works on rapidly approaching objects however is linked to spacious jumps or steps \cite{yang_2026,zhao_2026}.
Kohlbrenner et al. \cite{kohlbrenner_2026} presents a learning-based solution, that shows the usage of proximity sensors for a full tracking to motion execution approach.
Their controller manages to keep the H1 at the initial position for some of the evasion motions.

\paragraph{Comparability}
Another topic is the lack comparability between the above mentioned humanoid approaching object avoidance controllers.
As most of those are build evaluating a full controller package it is difficult to evaluate the capabilities of the motion planning algorithm.
Yang et al. \cite{yang_2026} test their PAC-MAN method by launching balls from $\SI{2.5 \pm 0.5}{\metre}$ distant frontal $\SI{\pm 25}{\deg}$ cone.
They also report that some were falling towards the legs, while most hit torso or head to trigger duck as well as lean reactions.
This results in untested scenarios for shots from the back of the humanoid.
It is also unknown to the reader at which initial velocity projectiles where launched.
\\
For Kohlbrenner et al. \cite{kohlbrenner_2026} balls are spawned with a velocity of $\SI{5 \pm 1}{\frac{\metre}{\second}}$.
Approaching objects are sampled uniformly distances between $\SI{4}{\metre}$ and $\SI{6}{\metre}$ at a height of $\SI{2.5 \pm 0.5}{\metre}$ of the ground.
The ball has a diameter of $\SI{15}{\cm}$.
\\
These differences in general test setups make it difficult to compare the resulting success rates.

\paragraph{Problem Formulation}
\begin{itemize}
  \item bipedal stance
  \item alternative to impact resistance (push recovery)
  \item comparable and reproducible tests
\end{itemize}
\paragraph{Overview}

% We present an optimal collision-avoidance trajectory generator based on a differentiable collision model using \acrfull{sdf} functions.
% In our scenario, the H1-2 from Unitree is placed in a bipedal stance.
% As a result, the humanoid performs object avoidance motions that resemble the ones from the protagonist of the "The Matrix" movie.
% Project Neo can be leveraged in control approaches like \acrfull{mpc}, as well as to precompute feasible object avoidance motions.
% Additionally, we provide a workflow to test and quantify such an avoidance controller.
% It defines a reproducible setup for generating large numbers of dodge scenarios.
% \\
% A general overview of Project Neo is visualized in \cref{fig:intro:graphicalAbstract}.
% \begin{figure}[t]
%   \centering
%   \includegraphics[width=0.9\textwidth]{figures/01_introduction/graphicalAbstract.png}
%   \caption{
%   Project Neo consists of two major contributions: a motion generator for optimal object avoidance and a benchmark to thoroughly test it across a wide range of object trajectories.
% }
%   \label{fig:intro:graphicalAbstract}
% \end{figure}
%

